\documentclass[10pt,a4paper]{amsart}
\usepackage{amsmath,amssymb,amsthm,graphicx}
% For commutative diagrams we use Xy-pic
\usepackage[all]{xy}

% We use 2cell for 2-commutative diagrams.
\xyoption{2cell}
% Theorem environments.
%
\theoremstyle{plain}
\newtheorem{theorem}[subsection]{Theorem}
\newtheorem{proposition}[subsection]{Proposition}
\newtheorem{lemma}[subsection]{Lemma}

\theoremstyle{definition}
\newtheorem{definition}[subsection]{Definition}
\newtheorem{example}[subsection]{Example}
\newtheorem{exercise}[subsection]{Exercise}
\newtheorem{situation}[subsection]{Situation}

\theoremstyle{remark}
\newtheorem{remark}[subsection]{Remark}
\newtheorem{remarks}[subsection]{Remarks}

\numberwithin{equation}{subsection}

% Macros
%
\def\lim{\mathop{\mathrm{lim}}\nolimits}
\def\colim{\mathop{\mathrm{colim}}\nolimits}
\def\Spec{\mathop{\mathrm{Spec}}}
\def\Hom{\mathop{\mathrm{Hom}}\nolimits}
\def\Ext{\mathop{\mathrm{Ext}}\nolimits}
\def\SheafHom{\mathop{\mathcal{H}\!\mathit{om}}\nolimits}
\def\SheafExt{\mathop{\mathcal{E}\!\mathit{xt}}\nolimits}
\def\Sch{\mathit{Sch}}
\def\Mor{\mathop{\mathrm{Mor}}\nolimits}
\def\Ob{\mathop{\mathrm{Ob}}\nolimits}
\def\Sh{\mathop{\mathit{Sh}}\nolimits}
\def\NL{\mathop{N\!L}\nolimits}
\def\CH{\mathop{\mathrm{CH}}\nolimits}
\def\proetale{{pro\text{-}\acute{e}tale}}
\def\etale{{\acute{e}tale}}
\def\QCoh{\mathit{QCoh}}
\def\Ker{\mathop{\mathrm{Ker}}}
\def\Im{\mathop{\mathrm{Im}}}
\def\Coker{\mathop{\mathrm{Coker}}}
\def\Coim{\mathop{\mathrm{Coim}}}

% Boxtimes
%
\DeclareMathSymbol{\boxtimes}{\mathbin}{AMSa}{"02}

%
% Macros for moduli stacks/spaces
%
\def\QCohstack{\mathcal{QC}\!\mathit{oh}}
\def\Cohstack{\mathcal{C}\!\mathit{oh}}
\def\Spacesstack{\mathcal{S}\!\mathit{paces}}
\def\Quotfunctor{\mathrm{Quot}}
\def\Hilbfunctor{\mathrm{Hilb}}
\def\Curvesstack{\mathcal{C}\!\mathit{urves}}
\def\Polarizedstack{\mathcal{P}\!\mathit{olarized}}
\def\Complexesstack{\mathcal{C}\!\mathit{omplexes}}
% \Pic is the operator that assigns to X its picard group, usage \Pic(X)
% \Picardstack_{X/B} denotes the Picard stack of X over B
% \Picardfunctor_{X/B} denotes the Picard functor of X over B
\def\Pic{\mathop{\mathrm{Pic}}\nolimits}
\def\Picardstack{\mathcal{P}\!\mathit{ic}}
\def\Picardfunctor{\mathrm{Pic}}
\def\Deformationcategory{\mathcal{D}\!\mathit{ef}}
\usepackage{cleveref}
\crefname{originalreference}{original reference}{original references}
\Crefname{originalreference}{Original reference}{Original references}
\makeatletter
\newlabel{remark-base-change-holds}{{0F02}{source}{}{sourcepointer-1}{}}
\newlabel{remark-base-change-holds@cref}{{[originalreference][0][]0F02}{[1][0][]source}{}{}{}}
\makeatother
\title{From BC(f, n, q-1) one deduces injectivity of the q-th base change map and a description of its cokernel}
\author{The Stacks Project Authors}
\date{stacks-project:etale-cohomology.tex:local-snapshot-20261011; GNU FDL 1.2 or later}
\begin{document}
\maketitle
\noindent\textit{Editorial note.} All delivered mathematical text is the author's own, copied from the pinned Stacks Project chapter file etale-cohomology.tex; no statement or proof step was rewritten, reordered or removed.\par\medskip
\noindent\textit{Editorial note.} The counted proof is the author's complete argument: the proof is written out in full: the choice of the injective embedding, the diagram of the two long exact sequences with the vanishing in the upper right corner, the use of BC(f, n, q-1) for the left vertical arrows, and the refinement of the embedding for part (2).\par\medskip
\noindent\textit{Editorial note.} The label lemma-base-change-q-injective names the counted statement, and the proof environment immediately following it in the source is delivered verbatim as its proof.\par\medskip
\noindent\textit{Editorial note.} The counted statement is carried by the author's own lemma environment in the chapter source.\par\medskip
\noindent\textit{Editorial note.} The statement is delivered verbatim with its displayed diagram and the three enumerated conclusions; the single cited input is kept as an external pointer with its canonical Stacks tag.\par\medskip
\noindent\textit{Editorial note.} The preamble is the chapter's own native preamble excerpt already used for this source: the Xy-pic setup (native preamble lines 14-18) and the native theorem declarations with the mathematical macros (native preamble lines 154-217).\par\medskip
\noindent\textit{Allowed standard prerequisites (editorial).}\begin{itemize}
\item The base change condition BC(f, n, q) and the standing remark. Same chapter, the remark collecting the base change conditions for the fixed morphism f and integer n, with the diagram of schemes used in the statement. Supplies the hypothesis BC(f, n, q-1) and the notation of the diagram.
\item Long exact cohomology sequences for the two rows. Same chapter, the long exact sequences of the two rows obtained from a short exact sequence of sheaves annihilated by n. Supplies the exact rows of the displayed diagram, including the vanishing in the upper right corner for an injective I.
\end{itemize}\medskip
\ifdefined\hypertarget\hypertarget{sourcepointer-1}{}\fi
\noindent\textit{Source reference.} Original reference 0F02: sources/stacks-etale-cohomology/pinned-original-source.tex, source lines 16848--16848. Prerequisite pointer invoked at proof line(s) 16870; it names a standard fact of the same chapter that the counted argument uses, and every step of the counted proof built on it is retained verbatim in the delivered body.\par\medskip
\begin{lemma}
\label{lemma-base-change-q-injective}
With $f : X \to S$ and $n$ as in Remark \ref{remark-base-change-holds}
assume for some $q \geq 1$ we have $BC(f, n, q - 1)$. Then
for every commutative diagram
$$
\xymatrix{
X \ar[d]_f & X' \ar[l] \ar[d]_{f'} & Y \ar[l]^h \ar[d]^e \\
S & S' \ar[l] & T \ar[l]_g
}
$$
with $X' = X \times_S S'$ and $Y = X' \times_{S'} T$ and
$g$ quasi-compact and quasi-separated, and every abelian sheaf
$\mathcal{F}$ on $T_\etale$ annihilated by $n$
\begin{enumerate}
\item the base change map
$(f')^{-1}R^qg_*\mathcal{F}\to R^qh_*e^{-1}\mathcal{F}$
is injective,
\item if $\mathcal{F} \subset \mathcal{G}$ where $\mathcal{G}$
on $T_\etale$ is annihilated by $n$, then
$$
\Coker\left(
(f')^{-1}R^qg_*\mathcal{F}\to R^qh_*e^{-1}\mathcal{F}
\right)
\subset
\Coker\left(
(f')^{-1}R^qg_*\mathcal{G}\to R^qh_*e^{-1}\mathcal{G}
\right)
$$
\item if in (2) the sheaf $\mathcal{G}$ is an injective sheaf
of $\mathbf{Z}/n\mathbf{Z}$-modules, then
$$
\Coker\left((f')^{-1}R^qg_*\mathcal{F}\to R^qh_*e^{-1}\mathcal{F} \right)
\subset R^qh_*e^{-1}\mathcal{G}
$$
\end{enumerate}
\end{lemma}

\begin{proof}
Choose a short exact sequence
$0 \to \mathcal{F} \to \mathcal{I} \to \mathcal{Q} \to 0$
where $\mathcal{I}$ is an injective sheaf of $\mathbf{Z}/n\mathbf{Z}$-modules.
Consider the induced diagram
$$
\xymatrix{
(f')^{-1}R^{q - 1}g_*\mathcal{I} \ar[d]_{\cong} \ar[r] &
(f')^{-1}R^{q - 1}g_*\mathcal{Q} \ar[d]_{\cong} \ar[r] &
(f')^{-1}R^qg_*\mathcal{F} \ar[d] \ar[r] &
0 \ar[d] \\
R^{q - 1}h_*e^{-1}\mathcal{I} \ar[r] &
R^{q - 1}h_*e^{-1}\mathcal{Q} \ar[r] &
R^qh_*e^{-1}\mathcal{F} \ar[r] &
R^qh_*e^{-1}\mathcal{I}
}
$$
with exact rows. We have the zero in the right upper corner
as $\mathcal{I}$ is injective. The left two vertical arrows are
isomorphisms by $BC(f, n, q - 1)$. We conclude that part (1) holds.
The above also shows that
$$
\Coker\left(
(f')^{-1}R^qg_*\mathcal{F}\to R^qh_*e^{-1}\mathcal{F}
\right)
\subset
R^qh_*e^{-1}\mathcal{I}
$$
hence part (3) holds. To prove (2) choose
$\mathcal{F} \subset \mathcal{G} \subset \mathcal{I}$.
\end{proof}

\end{document}
